\documentclass[11pt]{article}
\usepackage[letterpaper,margin=1in]{geometry}
\usepackage{amsmath,enumitem,xcolor}
\usepackage[hidelinks]{hyperref}
\setlength{\parskip}{0.35em}\setlength{\parindent}{0pt}
\setlist{itemsep=0.1em,topsep=0.2em,leftmargin=1.6em}
\definecolor{place}{RGB}{176,58,46}
\definecolor{open}{RGB}{31,97,141}
% where a figure or table goes, on its own line
\newcommand{\Place}[2]{\par\smallskip\noindent\fbox{\parbox{0.97\linewidth}{%
  \textcolor{place}{\textbf{#1}}\quad #2}}\par\smallskip}
% a decision still to be made
\newcommand{\Open}[1]{\textcolor{open}{\textit{Open: #1}}}
\title{Manuscript outline (reworked framing)}
\author{enceladus\_DSM \quad \texttt{studies/keff\_sintering/manuscript\_outline.tex}}
\date{2026-10-08}
\begin{document}
\maketitle

\textbf{Framing.} Sintering of ice by vapor transport, and the bulk
conductivity that results, under conditions a laboratory can reach; Enceladus
is the motivation and the closing implication. Every claim stays inside the
simulated range: $-40$ to $-5\,^{\circ}$C, air at 1~atm, sealed pores,
50~$\mu$m grains, 30~days, two dimensions. Themes below are prompts, not text.
Figures are in \texttt{ManuscriptFigures/}; tables are commands in
\texttt{Tables/tables.tex}. Notes behind this outline:
\texttt{MANUSCRIPT\_PLAN.md}.

\section*{Front matter}
\begin{itemize}
\item \textbf{Key points} (three). (i) For a given packing, temperature only
  rescales time. (ii) Conductivity follows specific surface area by one power
  law in well-connected packings; porosity sets the level. (iii) The model is
  the vapor route alone; what it does and does not reproduce of laboratory neck
  growth.
\item \textbf{Abstract}; \textbf{Plain Language Summary}. Written last.
\end{itemize}

\section{Introduction}
\begin{itemize}
\item \textbf{P1. The problem.} Ice aggregates change their thermal and
  mechanical properties as grains sinter; property evolution has been measured
  in the laboratory but not tied to microstructure.
\item \textbf{P2. Mechanisms.} Sintering without densification; vapor
  transport (sublimation and desublimation) and surface diffusion. Define
  ``desublimation'' here and reserve ``deposit'' for plume material.
\Place{Figure 1}{Mechanisms on a grain pair (a); the same process in an
  aggregate (b). Cite at the end of P2.}
\item \textbf{P3. What laboratory work has shown.} Molaro et al.\ (2019):
  neck growth, mechanisms by temperature. Choukroun et al.\ (2020):
  strengthening of sealed ice-powder samples, activation energy 24~kJ\,mol$^{-1}$.
  Demmenie et al.\ (2025): neck growth at saturation.
\item \textbf{P4. Why not natural snow.} Angular crystals at far higher
  porosity; overburden and densification. The analogues are unloaded
  ice-sphere and ice-powder samples.
\item \textbf{P5. Motivation.} Plume deposits on Enceladus: what a lander or
  a thermal measurement would need to know.
\item \textbf{P6. This work.} Phase-field model of vapor-transport sintering,
  periodic homogenization for conductivity, a porosity--temperature matrix;
  the three results; road map.
\end{itemize}

\section{Methods}
\subsection{Phase-field model of sintering by vapor transport}
\begin{itemize}
\item Fields ($\phi$, $T$, $\rho_v$) and governing equations; what is and is
  not in the model (no surface diffusion, no densification).
\item Material properties: cite \textbf{Table S1} (\verb|\TableMaterialProperties|).
\end{itemize}
\subsection{Interface kinetics and the Gibbs--Thomson condition}
\begin{itemize}
\item $\alpha_c$ and $\beta_{\mathrm{sub}}$ (always quoted together), $d_0$,
  $\rho_{vs}^{I}(T)$. Values: $\alpha_c = 1.020 \times 10^{-3}$ in the aggregates,
  $0.1$ in the Molaro pairs, $1.002 \times 10^{-3}$ in the saturated pair.
\item The thin-interface term $\Delta\beta$ in $\tau_{\mathrm{sub}}$, as in Kaempfer
  \& Plapp (2009) and Moure \& Fu (2024); what it costs (realised coefficient
  $\beta_{\mathrm{sub}} + \Delta\beta$).
\item The time scale $\tau_{\mathrm{sub}}$ and the sintering age
  $\theta = t/\tau_{\mathrm{sub}}$; state that $\varepsilon$ is fixed across
  the study. \Open{state the clock as $\tau_{\mathrm{sub}}$ at fixed
  $\varepsilon$, or as $\tau_R = \bar R^{2}\beta_{\mathrm{sub}}/d_0$.}
\end{itemize}
\subsection{Homogenization for the effective thermal conductivity}
\begin{itemize}
\item Cell problem, correctors, tensor conductivity law in the interface.
\Place{Figure 2}{The method on one packing: cell, interface and mesh,
  corrector, heat flux. After the cell problem is stated.}
\end{itemize}
\subsection{Numerical implementation}
\begin{itemize}
\item Discretization, time stepping, sampling of $\mathbf{K}_{\mathrm{eff}}$.
  Cite \textbf{Table S2} (\verb|\TableKeffParametersFull|).
\item Checks, each one sentence with a pointer to the supplement: domain size
  (Fig.~S1), interface width (Fig.~S2), time step and sampling (Fig.~S3).
\end{itemize}
\subsection{Aggregates: packings, boundary conditions, simulation matrix}
\begin{itemize}
\item Drop-and-roll deposition with void filling; five packings per porosity.
\item Periodic boundaries: scale separation, no wall physics to invent, a
  closed system driven by curvature alone.
\item The matrix: five porosities $\times$ five temperatures $\times$ five
  packings (125 runs). The master simulation used for single-run figures.
\end{itemize}

\section{Results}
\subsection{Two-grain sintering}
\textbf{Unsaturated pair (Molaro et al., 2019).}
\begin{itemize}
\item Set-up; cite \textbf{Table S3} (\verb|\TableGrainPairParameters|).
\item With the vapor loss matched to the measured shrinkage (a calibration at
  $-20\,^{\circ}$C), the model gives 41\% of the measured neck growth.
\item $-5\,^{\circ}$C: the authors' own caveats on that data (Sect.~4.1,
  Appendix~B).
\Place{Figure 3}{Sections, neck width, grain shrinkage. After the set-up
  paragraph.}
\end{itemize}
\textbf{Saturated pair (conditions of Demmenie et al., 2025).}
\begin{itemize}
\item Equal grains, vapor saturated over the grains' own curvature; the grain
  size holds.
\item Fitted exponent 0.21 after relaxation (0.21--0.23 by window), against
  0.26--0.33 measured and 1/3 for attachment-limited vapor transport.
  \Open{interpretation. Parallel mechanisms add rates, so a missing
  surface-diffusion term cannot raise an exponent; candidates are the mixed
  attachment/diffusion regime, the large initial neck, viscous flow of the
  quasi-liquid layer at $-3\,^{\circ}$C.}
\Place{Figure 4}{Sections and neck width with the free and one-third fits.
  After the result is stated.}
\end{itemize}

\subsection{Aggregates}
\textbf{Set-up.}
\begin{itemize}
\item One paragraph: the matrix and what is held fixed.
\Place{Table 1}{\texttt{\char`\\TableKeffParameters}. After the sentence
  that states the matrix.}
\end{itemize}
\textbf{What sintering does to the aggregate.}
\begin{itemize}
\item Necks form, small pores round and close, the ice fraction is conserved;
  orientation only, no analysis.
\Place{Figure 5}{Gallery: five porosities at four instants.}
\end{itemize}
\textbf{Temperature sets the clock.}
\begin{itemize}
\item Define $\theta$ in use; 30~days spans 43 to 1256~$\tau_{\mathrm{sub}}$.
\Place{Table 2}{\texttt{\char`\\TableKeffTemperature}. After $\theta$
  is introduced, before Figure 6.}
\item Five temperatures give one curve for $k_{\mathrm{eff}}$ and for SSA;
  for every packing $k_{\mathrm{eff}}$ at matched SSA agrees within 0.27\%
  (Fig.~S4 for all packings).
\item SSA falls linearly in $\log\theta$ after the transient, about 0.13 per
  decade; state the range, do not extrapolate.
\Place{Figure 6}{Snapshots of the master packing; $k_{\mathrm{eff}}$ against
  days and against $\theta$; SSA against $\theta$.}
\end{itemize}
\textbf{Conductivity as a function of the microstructural state.}
\begin{itemize}
\item Porosity dominates the level (a factor of about 3.3 across the range at
  30~d); packing-to-packing scatter is comparable to the gap between
  neighbouring porosities.
\item One power law in SSA, exponent $-0.80$, for porosity $\le 0.375$; the
  level falls exponentially with porosity, 6.6\% scatter.
\item Above 0.40 the exponent weakens and scatters; no claim made there.
\item Rise table by porosity and temperature (in text or Fig.~S5); the two
  densest packings tie.
\Place{Figure 7}{$k_{\mathrm{eff}}(t)$ by porosity; $k_{\mathrm{eff}}$ against
  SSA; level; exponent.}
\end{itemize}

\section{Discussion}
\subsection{Temperature dependence of the kinetics}
\begin{itemize}
\item Why the collapse holds: attachment-limited transfer, fixed $\alpha_c$,
  one discretization; what would break it.
\item Activation energy of the vapor route, about 50~kJ\,mol$^{-1}$, against
  24 measured by Choukroun et al.\ for strengthening; the ranges overlap at
  233--243~K, so this is a test an experiment can apply.
\end{itemize}
\subsection{SSA as the state variable}
\begin{itemize}
\item What the power law means and where it comes from; why porosity alone is
  not enough (packing scatter, connectivity).
\item Loss of the law at high porosity: poorly connected ice, amplified in two
  dimensions. Anisotropy as a consistency check (one sentence).
\end{itemize}
\subsection{What the vapor route captures and what it lacks}
\begin{itemize}
\item Read Figures 3 and 4 together: part of the measured neck growth, a lower
  exponent at saturation.
\item Surface diffusion and the quasi-liquid layer are absent; what each would
  change (rate versus exponent). Literature review needed before any claim.
\end{itemize}
\subsection{What a laboratory could measure}
\begin{itemize}
\item The three predictions as experiments: time--temperature superposition
  and its activation energy; $k_{\mathrm{eff}}$ against SSA on unloaded
  ice-sphere samples; SSA against log time.
\end{itemize}
\subsection{Implications for plume deposits on Enceladus}
\begin{itemize}
\item Qualitative only: sealed near-subsurface (about a metre down) as the
  closest setting; which conclusions carry (state dependence, exponent) and
  which do not (rates at 60--180~K, absolute conductivity).
\item No timescales in years. Small-sensor point, one or two sentences.
\end{itemize}
\subsection{Limitations}
\begin{itemize}
\item Two dimensions; air-filled pores at 1~atm; one grain size; vapor route
  only; first day or so depends on the interface width; claims for porosity
  $\le 0.375$.
\end{itemize}

\section{Conclusions}
\begin{itemize}
\item Three results, each with its range of validity; the single most useful
  experiment to do next.
\end{itemize}

\section*{Open Research; Acknowledgments}
Code and data availability; run inventory (\texttt{RUN\_TABLE.md}).

\section*{Supporting Information}
\begin{itemize}
\item \textbf{Tables.} S1 material properties; S2 complete simulation
  parameters; S3 two-grain parameters.
\item \textbf{Figures.} S1 domain-size convergence (and $k_{xy}$ against
  size); S2 interface width; S3 time step and sampling; S4 temperature collapse
  for all packings; S5 porosity trends, SSA by porosity, anisotropy in time.
\item \textbf{Text.} Packing construction; run table.
\end{itemize}

\end{document}
